%% Revision reference list (src/make_bib_rev.py). APA form, in order of first
%% citation, so the first work cited is [1]. paper/references_revision.md
%% records how each added entry was verified.

\begin{thebibliography}{99}

\bibitem{moffatt2024} %
\textit{Moffatt v.\ Air Canada}, 2024 BCCRT 149 (British Columbia Civil
Resolution Tribunal, 14 February 2024).

\bibitem{magesh2025hallucinationfree} %
Magesh, V., Surani, F., Dahl, M., Suzgun, M., Manning, C. D., \& Ho, D. E.
(2025). Hallucination-free? Assessing the reliability of leading AI legal
research tools. \textit{Journal of Empirical Legal Studies}, \textit{22}(2),
216--242.

\bibitem{beauchemin2023meaningbert} %
Beauchemin, D., Saggion, H., \& Khoury, R. (2023). MeaningBERT: Assessing
meaning preservation between sentences. \textit{Frontiers in Artificial
Intelligence}, \textit{6}, 1223924.

\bibitem{beauchemin2025judgebert} %
Beauchemin, D., Albert-Rochette, M., Khoury, R., \& D\'eziel, P.-L. (2025).
JUDGEBERT: Assessing legal meaning preservation between sentences. In
\textit{Proceedings of the 2025 Conference on Empirical Methods in Natural
Language Processing} (pp.~92--118).

\bibitem{sulem2018bleu} %
Sulem, E., Abend, O., \& Rappoport, A. (2018). BLEU is not suitable for the
evaluation of text simplification. In \textit{Proceedings of the 2018 Conference
on Empirical Methods in Natural Language Processing} (pp.~738--744).

\bibitem{alvamanchego2021unsuitability} %
Alva-Manchego, F., Scarton, C., \& Specia, L. (2021). The (un)suitability of
automatic evaluation metrics for text simplification. \textit{Computational
Linguistics}, \textit{47}(4), 861--889. \url{https://doi.org/10.1162/coli_a_00418}

\bibitem{garimella2022legalsimp} %
Garimella, A., Sancheti, A., Aggarwal, V., Ganesh, A., Chhaya, N., \&
Kambhatla, N. (2022). Text simplification for legal domain: Insights and
challenges. In \textit{Proceedings of the Natural Legal Language Processing
Workshop 2022} (pp.~296--304).

\bibitem{gardner2020contrast} %
Gardner, M., Artzi, Y., Basmov, V., Berant, J., Bogin, B., Chen, S., Dasigi, P.,
et al. (2020). Evaluating models' local decision boundaries via contrast sets.
In \textit{Findings of the Association for Computational Linguistics: EMNLP
2020} (pp.~1307--1323).

\bibitem{ribeiro2020checklist} %
Ribeiro, M. T., Wu, T., Guestrin, C., \& Singh, S. (2020). Beyond accuracy:
Behavioral testing of NLP models with CheckList. In \textit{Proceedings of the
58th Annual Meeting of the Association for Computational Linguistics}
(pp.~4902--4912).

\bibitem{zhang2019paws} %
Zhang, Y., Baldridge, J., \& He, L. (2019). PAWS: Paraphrase adversaries from
word scrambling. In \textit{Proceedings of the 2019 Conference of the North
American Chapter of the Association for Computational Linguistics: Human
Language Technologies} (pp.~1298--1308).

\bibitem{naik2018stress} %
Naik, A., Ravichander, A., Sadeh, N., Rose, C., \& Neubig, G. (2018). Stress test
evaluation for natural language inference. In \textit{Proceedings of the 27th
International Conference on Computational Linguistics} (pp.~2340--2353).

\bibitem{hossain2020negation} %
Hossain, M. M., Kovatchev, V., Dutta, P., Kao, T., Wei, E., \& Blanco, E.
(2020). An analysis of natural language inference benchmarks through the lens
of negation. In \textit{Proceedings of the 2020 Conference on Empirical Methods
in Natural Language Processing} (pp.~9106--9118).

\bibitem{sai2021perturbation} %
Sai, A. B., Dixit, T., Sheth, D. Y., Mohan, S., \& Khapra, M. M. (2021).
Perturbation CheckLists for evaluating NLG evaluation metrics. In
\textit{Proceedings of the 2021 Conference on Empirical Methods in Natural
Language Processing} (pp.~7219--7234).

\bibitem{he2023blindspots} %
He, T., Zhang, J., Wang, T., Kumar, S., Cho, K., Glass, J., \& Tsvetkov, Y.
(2023). On the blind spots of model-based evaluation metrics for text
generation. In \textit{Proceedings of the 61st Annual Meeting of the
Association for Computational Linguistics (Volume 1: Long Papers)}
(pp.~12067--12097).

\bibitem{chen2023menli} %
Chen, Y., \& Eger, S. (2023). MENLI: Robust evaluation metrics from natural
language inference. \textit{Transactions of the Association for Computational
Linguistics}, \textit{11}, 804--825.

\bibitem{mujahid2026stress} %
Mujahid, Z. M., Wright, D., \& Augenstein, I. (2026). Stress testing factual
consistency metrics for long-document summarization. In \textit{Proceedings of
the 64th Annual Meeting of the Association for Computational Linguistics (Volume
1: Long Papers)} (pp.~31914--31933).

\bibitem{goebel2024coliee} %
Goebel, R., Kano, Y., Kim, M.-Y., Rabelo, J., Satoh, K., \& Yoshioka, M. (2024).
Overview and discussion of the Competition on Legal Information,
Extraction/Entailment (COLIEE) 2023. \textit{The Review of Socionetwork
Strategies}, \textit{18}(1), 27--47.

\bibitem{koreeda2021contractnli} %
Koreeda, Y., \& Manning, C. (2021). ContractNLI: A dataset for document-level
natural language inference for contracts. In \textit{Findings of the
Association for Computational Linguistics: EMNLP 2021} (pp.~1907--1919).

\bibitem{guha2023legalbench} %
Guha, N., Nyarko, J., Ho, D. E., R\'e, C., Chilton, A., Narayana, A.,
Chohlas-Wood, A., et al. (2023). LegalBench: A collaboratively built benchmark
for measuring legal reasoning in large language models. \textit{Advances in
Neural Information Processing Systems}, \textit{36} (Datasets and Benchmarks
Track).

\bibitem{zheng2023judging} %
Zheng, L., Chiang, W.-L., Sheng, Y., Zhuang, S., Wu, Z., Zhuang, Y., Lin, Z.,
et al. (2023). Judging LLM-as-a-judge with MT-Bench and Chatbot Arena.
\textit{Advances in Neural Information Processing Systems}, \textit{36}.

\bibitem{enguehard2025lemaj} %
Enguehard, J., Van Ermengem, M., Atkinson, K., Cha, S., Ghosh Chowdhury, A.,
Kallur Ramaswamy, P., Roghair, J., et al. (2025). LeMAJ (Legal LLM-as-a-Judge):
Bridging legal reasoning and LLM evaluation. In \textit{Proceedings of the
Natural Legal Language Processing Workshop 2025} (pp.~318--337).

\bibitem{bean2025construct} %
Bean, A. M., Kearns, R. O., Romanou, A., Hafner, F. S., Mayne, H., Batzner, J.,
Foroutan Eghlidi, N., et al. (2025). Measuring what matters: Construct validity
in large language model benchmarks. \textit{Advances in Neural Information
Processing Systems}, \textit{38}, 19868--19949.

\bibitem{pacchiardi2024cleverhans} %
Pacchiardi, L., Tesic, M., Cheke, L. G., \& Hern\'andez-Orallo, J. (2024).
Leaving the barn door open for Clever Hans: Simple features predict LLM
benchmark answers. \textit{arXiv preprint arXiv:2410.11672}.

\bibitem{deutsch2021statistical} %
Deutsch, D., Dror, R., \& Roth, D. (2021). A statistical analysis of
summarization evaluation metrics using resampling methods.
\textit{Transactions of the Association for Computational Linguistics},
\textit{9}, 1132--1146.

\bibitem{steiger1980tests} %
Steiger, J. H. (1980). Tests for comparing elements of a correlation matrix.
\textit{Psychological Bulletin}, \textit{87}(2), 245--251.

\bibitem{graham2014williams} %
Graham, Y., \& Baldwin, T. (2014). Testing for significance of increased
correlation with human judgment. In \textit{Proceedings of the 2014 Conference
on Empirical Methods in Natural Language Processing (EMNLP)} (pp.~172--176).

\bibitem{plank2022variation} %
Plank, B. (2022). The ``problem'' of human label variation: On ground truth in
data, modeling and evaluation. In \textit{Proceedings of the 2022 Conference on
Empirical Methods in Natural Language Processing} (pp.~10671--10682).

\bibitem{xu2026disagreement} %
Xu, Y., \& Jurgens, D. (2026). Beyond consensus: Perspectivist modeling and
evaluation of annotator disagreement in NLP. \textit{arXiv preprint
arXiv:2601.09065}.

\bibitem{yang2026compared} %
Yang, Z., Levy, M., Goldberg, Y., \& Wallace, B. C. (2026). Compared to what?
Baselines and metrics for counterfactual prompting. \textit{arXiv preprint
arXiv:2605.01048}.

\bibitem{interpretationact} %
\textit{Interpretation Act}, CQLR c I-16, s 51 (Quebec).

\bibitem{zhang2019bertscore} %
Zhang, T., Kishore, V., Wu, F., Weinberger, K. Q., \& Artzi, Y. (2020).
BERTScore: Evaluating text generation with BERT. In \textit{International
Conference on Learning Representations}.

\bibitem{laurer2024nli} %
Laurer, M., van Atteveldt, W., Casas, A., \& Welbers, K. (2024). Less
annotating, more classifying: Addressing the data scarcity issue of supervised
machine learning with deep transfer learning and BERT-NLI. \textit{Political
Analysis}, \textit{32}(1), 84--100.

\bibitem{uma2021learning} %
Uma, A. N., Fornaciari, T., Hovy, D., Paun, S., Plank, B., \& Poesio, M. (2021).
Learning from disagreement: A survey. \textit{Journal of Artificial Intelligence
Research}, \textit{72}, 1385--1470.

\bibitem{peterson2019human} %
Peterson, J. C., Battleday, R. M., Griffiths, T. L., \& Russakovsky, O. (2019).
Human uncertainty makes classification more robust. In \textit{Proceedings of
the IEEE/CVF International Conference on Computer Vision (ICCV)}
(pp.~9616--9625).

\bibitem{koenker1978regression} %
Koenker, R., \& Bassett, G., Jr. (1978). Regression quantiles.
\textit{Econometrica}, \textit{46}(1), 33--50.

\bibitem{krippendorff2004content} %
Krippendorff, K. (2019). \textit{Content analysis: An introduction to its
methodology}. SAGE. \url{https://doi.org/10.4135/9781071878781}

\bibitem{kaushik2020counterfactual} %
Kaushik, D., Hovy, E., \& Lipton, Z. C. (2020). Learning the difference that
makes a difference with counterfactually-augmented data. In
\textit{International Conference on Learning Representations}.

\bibitem{huang2020counterfactual} %
Huang, W., Liu, H., \& Bowman, S. R. (2020). Counterfactually-augmented SNLI
training data does not yield better generalization than unaugmented data. In
\textit{Proceedings of the First Workshop on Insights from Negative Results in
NLP} (pp.~82--87).

\end{thebibliography}
